\expandafter\gdef\csname SNside1\endcsname{\node[anchor=north west,inner sep=0pt,text width=48mm,font=\SideFont,align=left] at ([xshift=138mm,yshift=-565.45pt]current page.north west) {\textcolor{SNink}{\textbf{错项排除 · 同号不一定同文件}}\par\vspace{1.5mm}inode号要在所属文件系统内解释。两个不同分区上都出现inode10，不足以证明它们指向同一文件；同一文件系统中，两个名字解析到同一个inode才表示这种对象共享。判题时先确定命名空间，再比较数字。};}
\expandafter\gdef\csname SNside2\endcsname{\node[anchor=north west,inner sep=0pt,text width=48mm,font=\SideFont,align=left] at ([xshift=138mm,yshift=-301.97pt]current page.north west) {\textcolor{SNink}{\textbf{最小反例 · 关闭只撤销本次使用引用}}\par\vspace{1.5mm}一个文件有名字a，进程打开后再关闭，目录里的a仍存在，文件数据也应保留。close减少的是打开使用关系，unlink取消的是名字关系。只有满足没有硬链接、没有有效使用引用等条件，才讨论对象最终回收。};}
\expandafter\gdef\csname SNside3\endcsname{\node[anchor=north west,inner sep=0pt,text width=48mm,font=\SideFont,align=left] at ([xshift=138mm,yshift=-95.00pt]current page.north west) {\textcolor{SNink}{\textbf{不变量 · 文件身份不决定当前偏移}}\par\vspace{1.5mm}独立open同一inode可产生两个打开实例，分别维护位置；dup或继承则可能共享同一个打开实例及位置。始终沿``fd到打开实例到inode''画引用箭头，就能判断哪项状态共享，而不是仅凭文件名相同推结论。};}
\expandafter\gdef\csname SNside4\endcsname{\node[anchor=north west,inner sep=0pt,text width=48mm,font=\SideFont,align=left] at ([xshift=138mm,yshift=-327.83pt]current page.north west) {\textcolor{SNink}{\textbf{错项排除 · 权限集合不能压成一个编号}}\par\vspace{1.5mm}每类用户可能同时有读、写等多种权限；若有5个独立许可就要表达其全部子集，需要5位。给用户类别和某一操作编码，只能描述一次选择，不能保存每一类用户各自的完整许可组合，二者信息量不同。};}
\expandafter\gdef\csname SNside5\endcsname{\node[anchor=north west,inner sep=0pt,text width=48mm,font=\SideFont,align=left] at ([xshift=138mm,yshift=-95.00pt]current page.north west) {\textcolor{SNink}{\textbf{课程连接 · 两个索引属于不同层次}}\par\vspace{1.5mm}逻辑记录索引把关键字对应到记录或文件内偏移；物理索引把逻辑块号对应到磁盘块。数据结构中的查找复杂度可解释前一层，操作系统的间接索引层数解释后一层。遇到``索引''先写它的输入与输出对象。};}
\expandafter\gdef\csname SNside6\endcsname{\node[anchor=north west,inner sep=0pt,text width=48mm,font=\SideFont,align=left] at ([xshift=138mm,yshift=-95.00pt]current page.north west) {\textcolor{SNink}{\textbf{错项排除 · FAT少读数据块仍需追链}}\par\vspace{1.5mm}FAT驻留内存时，定位后面簇不必逐个读取前面数据块，但仍要沿内存中的链获取后继号。其优势是减少磁盘I/O，不能直接推成按逻辑下标常数时间定位；表规模还由整个受管卷的簇数决定。};}
\expandafter\gdef\csname SNside7\endcsname{\node[anchor=north west,inner sep=0pt,text width=48mm,font=\SideFont,align=left] at ([xshift=138mm,yshift=-156.65pt]current page.north west) {\textcolor{SNink}{\textbf{最小反例 · 删一个名字不等于删文件}}\par\vspace{1.5mm}同一文件inode为i7，目录项a和b都指向它。删除a后仍有b这一硬链接，文件内容继续可达。把a叫作原文件不赋予它特殊地位；硬链接之间对称，不能说删除``原名字''就必然使另一个硬链接失效。\par\smallskip\begin{tikzpicture}[x=1mm,y=1mm,>=Stealth,node distance=5mm]
\node[draw=SNink!50,fill=SNblue,rounded corners=1mm,text width=34mm,align=center,inner sep=1.2mm,font=\fontsize{7.5}{9}\selectfont] (m0) at (0,0) {a指向i7};
\node[draw=SNink!50,fill=SNblue,rounded corners=1mm,text width=34mm,align=center,inner sep=1.2mm,font=\fontsize{7.5}{9}\selectfont] (m1) at (0,-11) {b指向i7};
\node[draw=SNink!50,fill=SNblue,rounded corners=1mm,text width=34mm,align=center,inner sep=1.2mm,font=\fontsize{7.5}{9}\selectfont] (m2) at (0,-22) {删a留b};
\draw[->,SNink] (m0) edge[bend left=40] (m2);
\draw[->,SNink] (m1) -- (m2);
\end{tikzpicture}
\par{\fontsize{7.5}{9}\selectfont 两个目录项对应同一inode；删除名字a后，名字b仍可引用文件。}};}
\expandafter\gdef\csname SNside8\endcsname{\node[anchor=north west,inner sep=0pt,text width=48mm,font=\SideFont,align=left] at ([xshift=138mm,yshift=-95.00pt]current page.north west) {\textcolor{SNink}{\textbf{不变量 · 缓存只删去对应访问}}\par\vspace{1.5mm}目录内容已在内存，只免去读取该目录内容；若目标inode和数据尚未缓存，仍要分别取到。写访盘路径时为每个节点标``已缓存/未缓存''，再计边上的实际读取，避免一个命中条件被扩张成整条路径全命中。};}
\expandafter\gdef\csname SNside9\endcsname{\node[anchor=north west,inner sep=0pt,text width=48mm,font=\SideFont,align=left] at ([xshift=138mm,yshift=-95.00pt]current page.north west) {\textcolor{SNink}{\textbf{错项排除 · 位图长度取决于管理范围}}\par\vspace{1.5mm}管理一百万个块，即使当前只剩一个空闲块，位图仍须为每个块保留一位。空闲块数量影响位值分布，不缩短受管对象集合。位图内相对位置算完后再加其起始磁盘块号，不能先从受管块号减去起点。};}
\expandafter\gdef\csname SNside10\endcsname{\node[anchor=north west,inner sep=0pt,text width=48mm,font=\SideFont,align=left] at ([xshift=138mm,yshift=-95.00pt]current page.north west) {\textcolor{SNink}{\textbf{课程连接 · 统一接口可以隐藏多种布局}}\par\vspace{1.5mm}VFS把查找、打开、读写等操作抽象成共同接口，再由具体文件系统实现。它和面向接口编程相似：调用约定统一，并不要求内部数据结构一致。不能因为应用都调用read，就推断各磁盘格式具有相同的inode字段。};}
\expandafter\gdef\csname SNside11\endcsname{\node[anchor=north west,inner sep=0pt,text width=48mm,font=\SideFont,align=left] at ([xshift=138mm,yshift=-95.00pt]current page.north west) {\textcolor{SNink}{\textbf{不变量 · 文件数量由最紧资源约束}}\par\vspace{1.5mm}每个图像既需要一个inode，也需要足够数据簇。应分别计算inode区能容纳的个数和数据区能容纳的个数，再取较小者；若还要间接索引簇或目录空间，应按题设另计，不能只用剩余字节除文件长度。};}
\expandafter\gdef\csname SNside12\endcsname{\node[anchor=north west,inner sep=0pt,text width=48mm,font=\SideFont,align=left] at ([xshift=138mm,yshift=-95.00pt]current page.north west) {\textcolor{SNink}{\textbf{不变量 · 文件数量由最紧资源约束}}\par\vspace{1.5mm}每个图像既需要一个inode，也需要足够数据簇。应分别计算inode区能容纳的个数和数据区能容纳的个数，再取较小者；若还要间接索引簇或目录空间，应按题设另计，不能只用剩余字节除文件长度。};}
\expandafter\gdef\csname SNside13\endcsname{\node[anchor=north west,inner sep=0pt,text width=48mm,font=\SideFont,align=left] at ([xshift=138mm,yshift=-95.00pt]current page.north west) {\textcolor{SNink}{\textbf{不变量 · 容量与索引开销分开累计}}\par\vspace{1.5mm}一个索引块可装B除以指针字节数个地址。计算最大文件长度时只累加可指向的数据块，索引块本身另占外存。若题目转问总空间消耗，才把实际使用的各级索引块加回，两个问题不能共用一个总数。};}
\expandafter\gdef\csname SNside14\endcsname{\node[anchor=north west,inner sep=0pt,text width=48mm,font=\SideFont,align=left] at ([xshift=138mm,yshift=-95.00pt]current page.north west) {\textcolor{SNink}{\textbf{错项排除 · 位图长度取决于管理范围}}\par\vspace{1.5mm}管理一百万个块，即使当前只剩一个空闲块，位图仍须为每个块保留一位。空闲块数量影响位值分布，不缩短受管对象集合。位图内相对位置算完后再加其起始磁盘块号，不能先从受管块号减去起点。};}
\AddToShipoutPictureFG{\begin{tikzpicture}[remember picture,overlay]\ifcsname SNside\thepage\endcsname\csname SNside\thepage\endcsname\fi\end{tikzpicture}}